\section{Geometry of finite legal intervals}\label{sec:silent-geometry}
The local language controls the geometry of an entire legal interval, even
when its duration is unbounded. The coordinate selected in the next theorem
is fixed for the whole interval.

\textbf{Theorem (endpoint strip).}
Let $p_0,\ldots,p_L$ be the head positions of a physically feasible finite
word whose contained ten-factors are standard. There is an $i\in\{x,y\}$ such that
\[
 \min(p_{0,i},p_{L,i})-43\le p_{t,i}\le
 \max(p_{0,i},p_{L,i})+43\qquad(0\le t\le L).
\]
More precisely, after translating $p_0$ to zero, one of the four coordinate
normals $n$ satisfies $-43\le n\cdot p_t\le n\cdot p_L+43$ throughout.

\textit{Proof.}
Use the bounded-fold lift established for the event macros. Adjoining the
incoming heading gives a directed graph $G$ with 30,380 vertices and 31,052
edges. A vertex records the length-149 history, the fold count and the actual
heading; an edge appends its literal turn and carries its unit displacement.
Every feasible continuation after the first 149 turns lifts to this graph.
The all-parameter exclusion of $BA^nBA^mB$ supplies coverage of the fold count.

For a cyclic strongly connected component $C$, let
$M(C)=\{n:n\cdot d(C)\ge0\}$, where $n$ ranges over the four coordinate
normals and $d(C)$ is the displacement of its simple cycle. Give an acyclic
component all four normals. The finite certificate starts a mask at every
component and propagates every possible intersection along every condensation
edge. All 88,024 resulting component/mask states have nonempty masks.
Closure under propagation proves by induction that every finite path has a
common allowed normal. An independent work-list calculation on the original
graph agrees at all 90,256 vertex/mask states.

For each normal delete the cyclic components that omit it. On the resulting
$G_n$, the certificate gives an integer potential $H_n$ with
\[
 0\le H_n\le33,\qquad
 n\cdot\Delta(v,u)\ge H_n(u)-H_n(v)
 \quad\text{for every edge }v\longrightarrow u.
\]
There are 30,714 inequalities for each normal. Telescoping bounds every
subpath's retreat by 33. These inequalities and the nonempty-mask closure,
rather than the method used to find the potentials, are the finite premises.

Short words of length at most 149 have coordinate radius at most 24 by
complete feasible-word enumeration. For a longer word put $q=p_{149}$,
$v$ equal to its startup lift state, and $z_t=n\cdot p_t$. The startup
certificate checks, for every such state and every allowed normal,
\[
 \min_{t\le149}z_t\ge-43,\quad z_{149}-H_n(v)\ge-43,\quad
 \max_{t\le149}(z_t-z_{149}+H_n(v))\le43.
\]
For $t\ge149$, telescoping gives $z_t\ge z_{149}-H_n(v)\ge-43$
and $z_t\le z_L+33$. For $t\le149$, use the first startup inequality
and $z_L\ge z_{149}-H_n(v)$ in the third. Short words use the sign
with $n\cdot p_L\ge0$. This proves the oriented assertion and the strip.
\hfill$\square$

\textbf{Corollary (localized events).}
Let $K$ be a nonempty finite set with bounding rectangle
$[x_-,x_+]\times[y_-,y_+]$. Between any consecutive chain of intrinsic
$J_{10}$ events with read positions in $K$, all head positions lie in
\[
 \{x_- -44\le x\le x_+ +44\}\ \cup\
 \{y_- -44\le y\le y_+ +44\}.
\]
The conclusion includes the outgoing move of the last event and is independent
of the number of events in the chain. Indeed, each intervening silent word
starts one step from $K$ and ends in $K$; apply the theorem to that contiguous
word and include the event moves. If infinitely many events are eventually
localized in $K$, the entire late head trajectory lies in this same cross.

This geometry complements the complete-source bounds used by event prediction.
A legal interval occupies one endpoint strip while it may travel arbitrarily
far along it. The folded ribbon in Appendix~\ref{app:folded-ribbon} realizes precisely this distinction.
